\chapter{Annales — Sujets complets type BEPC}

\noindent\textit{Ce chapitre propose trois épreuves complètes de Physique-Chimie-Technologie,
construites pour reproduire fidèlement la structure, la durée, le coefficient et le barème
réels de l'épreuve du BEPC (Partie A — Évaluation des ressources : savoirs et savoir-faire ;
Partie B — Évaluation des compétences), telle que décrite au chapitre « Méthodologie de
l'épreuve ». Contrairement aux annales officielles d'une session précise, ces trois sujets
sont des créations originales du fascicule, chacun croisant plusieurs chapitres du
programme autour d'une situation concrète, avec un corrigé entièrement rédigé et vérifié.}
\vspace{8pt}

\connecteBanner

% ═══════════════════════════════════════════════════════════════
\section{Sujet 1 — Le forage villageois}

\noindent\textbf{Examen :} BEPC \quad\textbf{Épreuve :} Physique-Chimie-Technologie
\quad\textbf{Durée :} 2 heures \quad\textbf{Coefficient :} 2 \quad\textbf{Notée sur :}
20 points.
\par\smallskip\noindent
L'épreuve comporte deux parties A et B, indépendantes et obligatoires.

\subsection*{Partie A — Évaluation des ressources (10 points)}

\paragraph{a) Évaluation des savoirs (5 points)}
\begin{enumerate}[label=\arabic*)]
\item Définir le mot « molécule ». \hfill (1pt)
\item Énoncer la loi de Lavoisier. \hfill (1pt)
\item Qu'est-ce qu'une machine simple ? \hfill (1pt)
\item Donner la relation entre la tension $U$, la résistance $R$ et l'intensité $I$ aux
bornes d'un conducteur ohmique (loi d'Ohm). \hfill (1pt)
\item Que signifie l'acronyme PMH, rencontré dans l'étude des moteurs à combustion
interne ? \hfill (1pt)
\end{enumerate}

\paragraph{b) Évaluation des savoir-faire (5 points)}
\begin{enumerate}[label=\arabic*)]
\item On dissout $n=0{,}6$ mol de sel dans un volume $V=3$ L d'eau. Calculer la
concentration molaire $c$ de la solution obtenue. \hfill (2pt)
\item Une pompe électrique fonctionne sous une tension $U=220$ V, traversée par un courant
d'intensité $I=5$ A. Calculer la puissance électrique de cette pompe. \hfill (1,5pt)
\item Un treuil de secours a un tambour de rayon $r=0{,}1$ m et une manivelle de longueur
$R=0{,}4$ m. Il sert à remonter un seau de poids $P=150$ N. Calculer la force $F$ à
exercer sur la manivelle. \hfill (1,5pt)
\end{enumerate}

\subsection*{Partie B — Évaluation des compétences (10 points)}

\noindent\textbf{Situation :} Le comité de développement du village de Ngola vient
d'installer un forage équipé d'une pompe électrique pour puiser l'eau. En cas de coupure de
courant, un treuil manuel permet de continuer à puiser l'eau à la main. Le comité doit
aussi désinfecter l'eau avant sa distribution, et souhaite estimer le coût de
fonctionnement de la pompe.

\begin{enumerate}[label=\textbf{Tâche \arabic*.}, leftmargin=2.2cm]
\item La pompe fonctionne sous une tension $U=220$ V, traversée par un courant
d'intensité $I=5$ A, pendant $t=3$ heures chaque jour. Calculer sa puissance électrique,
puis l'énergie qu'elle consomme chaque jour (en Wh, puis en kWh). \hfill (3pt)
\item Sachant que le kilowattheure coûte $79$ FCFA, calculer le coût journalier de
fonctionnement de la pompe. \hfill (2pt)
\item Lors d'une panne, le treuil de secours (tambour $r=0{,}15$ m, manivelle
$R=0{,}6$ m) est utilisé pour remonter un seau de poids $P=120$ N. Calculer la force $F$
à exercer. \hfill (2pt)
\item Pour désinfecter l'eau, on dissout $n=0{,}3$ mol de produit désinfectant dans un
bidon test de volume $V=1{,}5$ L. Calculer la concentration $c$ obtenue, en mol/L.
\hfill (2pt)
\item En cas de panne prolongée, le comité hésite entre le treuil manuel (tâche 3) et une
poulie mobile pour remonter la même charge de $120$ N. Lequel des deux dispositifs demande
le moins d'effort ? Conclure. \hfill (1pt)
\end{enumerate}

% ═══════════════════════════════════════════════════════════════
\section{Sujet 1 — Corrigé}

\subsection*{Partie A}

\paragraph{a) Savoirs}
\begin{enumerate}[label=\arabic*)]
\item Une molécule est un assemblage ordonné et électriquement neutre d'atomes liés entre
eux par des liaisons de covalence.
\item Loi de Lavoisier : au cours d'une réaction chimique, la masse totale des réactifs
consommés est égale à la masse totale des produits formés (la masse se conserve).
\item Une machine simple est un dispositif mécanique élémentaire qui permet de modifier
l'intensité ou la direction d'une force.
\item $U=R\times I$.
\item PMH signifie Point Mort Haut : la position la plus haute atteinte par le piston dans
le cylindre.
\end{enumerate}

\paragraph{b) Savoir-faire}
\begin{enumerate}[label=\arabic*)]
\item $c=\dfrac nV=\dfrac{0{,}6}3=0{,}2$ mol/L.
\item $P=U\times I=220\times5=1\,100$ W.
\item $F=\dfrac{P\times r}{R}=\dfrac{150\times0{,}1}{0{,}4}=37{,}5$ N.
\end{enumerate}

\subsection*{Partie B}

\textbf{Tâche 1.} $P=U\times I=220\times5=1\,100$ W. Énergie journalière :
$E=P\times t=1\,100\times3=3\,300$ Wh $=3{,}3$ kWh.

\textbf{Tâche 2.} Coût journalier $=3{,}3\times79=260{,}7$ FCFA.

\textbf{Tâche 3.} $F=\dfrac{P\times r}{R}=\dfrac{120\times0{,}15}{0{,}6}=30$ N.

\textbf{Tâche 4.} $c=\dfrac nV=\dfrac{0{,}3}{1{,}5}=0{,}2$ mol/L.

\textbf{Tâche 5.} Avec la poulie mobile : $F=\dfrac P2=\dfrac{120}2=60$ N, contre $30$ N
avec le treuil (tâche 3). Le treuil demande donc moins d'effort ($30$ N $<60$ N) : c'est le
dispositif le plus avantageux ici.

% ═══════════════════════════════════════════════════════════════
\section{Sujet 2 — L'atelier du garagiste}

\noindent\textbf{Examen :} BEPC \quad\textbf{Épreuve :} Physique-Chimie-Technologie
\quad\textbf{Durée :} 2 heures \quad\textbf{Coefficient :} 2 \quad\textbf{Notée sur :}
20 points.
\par\smallskip\noindent
L'épreuve comporte deux parties A et B, indépendantes et obligatoires.

\subsection*{Partie A — Évaluation des ressources (10 points)}

\paragraph{a) Évaluation des savoirs (5 points)}
\begin{enumerate}[label=\arabic*)]
\item Qu'est-ce qu'un corps simple ? Donner un exemple. \hfill (1pt)
\item Citer, dans l'ordre, les quatre temps du cycle d'un moteur à combustion interne.
\hfill (1pt)
\item Quelle est la différence entre un matériau thermoplastique et un matériau
thermodurcissable ? \hfill (1pt)
\item Que signifie une échelle $1{:}50$ indiquée sur un dessin technique ? \hfill (1pt)
\item Qu'est-ce qu'un système de transmission réducteur de vitesse ? \hfill (1pt)
\end{enumerate}

\paragraph{b) Évaluation des savoir-faire (5 points)}
\begin{enumerate}[label=\arabic*)]
\item Un moteur reçoit une énergie $E_{\text{fournie}}=250$ kJ de carburant et fournit une
énergie mécanique utile $E_{\text{utile}}=70$ kJ. Calculer son rendement $\eta$.
\hfill (2pt)
\item Une poulie menante de diamètre $D_1=60$ mm tourne à $N_1=2\,400$ tr/min et entraîne,
par une courroie, une poulie menée de diamètre $D_2=90$ mm. Calculer la vitesse de
rotation $N_2$ de la poulie menée. \hfill (2pt)
\item Une pièce mécanique réelle mesure $45$ cm de long. Elle est dessinée à l'échelle
$1{:}5$. Calculer sa longueur sur le dessin. \hfill (1pt)
\end{enumerate}

\subsection*{Partie B — Évaluation des compétences (10 points)}

\noindent\textbf{Situation :} Dans son atelier, un garagiste répare le moteur d'une
mobylette en panne et doit remplacer la courroie de transmission qui entraîne le
ventilateur de refroidissement. Il doit également dessiner la pièce défectueuse pour en
commander une neuve, et calcule sa consommation électrique pour l'outillage utilisé.

\begin{enumerate}[label=\textbf{Tâche \arabic*.}, leftmargin=2.2cm]
\item Le moteur de la mobylette reçoit une énergie de $E_{\text{fournie}}=250$ kJ par
cycle de test et restitue $E_{\text{utile}}=70$ kJ sous forme d'énergie mécanique.
Calculer son rendement $\eta$. \hfill (2pt)
\item La poulie du moteur, de diamètre $D_1=50$ mm, tourne à $N_1=3\,000$ tr/min et
entraîne, par une courroie, la poulie du ventilateur de diamètre $D_2=25$ mm. Calculer la
vitesse de rotation $N_2$ du ventilateur. \hfill (2pt)
\item Calculer le rapport de transmission $r$ de ce système. Est-il réducteur ou
multiplicateur de vitesse ? \hfill (2pt)
\item Le garagiste dessine la pièce défectueuse, qui mesure réellement $8$ cm de long, à
l'échelle $1{:}2$. Calculer sa longueur sur le dessin. \hfill (2pt)
\item Il utilise une perceuse électrique fonctionnant sous $U=220$ V, traversée par un
courant $I=2$ A. Calculer la puissance de cette perceuse, puis conclure sur l'intérêt
d'un moteur à haut rendement dans un atelier professionnel. \hfill (2pt)
\end{enumerate}

% ═══════════════════════════════════════════════════════════════
\section{Sujet 2 — Corrigé}

\subsection*{Partie A}

\paragraph{a) Savoirs}
\begin{enumerate}[label=\arabic*)]
\item Un corps simple est une substance dont les molécules ne contiennent qu'un seul type
d'atome. Exemple : le dioxygène $\mathrm{O_2}$.
\item Admission, compression, explosion (détente), échappement.
\item Le thermoplastique ramollit à la chaleur et peut être refondu ; le thermodurcissable
durcit définitivement et ne peut plus être refondu.
\item Une échelle $1{:}50$ signifie que chaque longueur sur le dessin représente $50$ fois
cette longueur dans la réalité (réduction).
\item Un système dans lequel la roue menée tourne moins vite que la roue menante
($r=\dfrac{N_2}{N_1}<1$), ce qui augmente la force transmise.
\end{enumerate}

\paragraph{b) Savoir-faire}
\begin{enumerate}[label=\arabic*)]
\item $\eta=\dfrac{70}{250}\times100=28\,\%$.
\item $N_2=N_1\times\dfrac{D_1}{D_2}=2\,400\times\dfrac{60}{90}=1\,600$ tr/min.
\item Dessin $=45\times\dfrac15=9$ cm.
\end{enumerate}

\subsection*{Partie B}

\textbf{Tâche 1.} $\eta=\dfrac{E_{\text{utile}}}{E_{\text{fournie}}}\times100=
\dfrac{70}{250}\times100=28\,\%$.

\textbf{Tâche 2.} $N_2=N_1\times\dfrac{D_1}{D_2}=3\,000\times\dfrac{50}{25}=6\,000$
tr/min.

\textbf{Tâche 3.} $r=\dfrac{N_2}{N_1}=\dfrac{6\,000}{3\,000}=2>1$ : le système est
multiplicateur de vitesse.

\textbf{Tâche 4.} Dessin $=8\times\dfrac12=4$ cm.

\textbf{Tâche 5.} $P=U\times I=220\times2=440$ W. Un moteur à rendement plus élevé
consomme moins de carburant pour un même travail fourni, ce qui réduit les coûts de
fonctionnement de l'atelier et son impact environnemental.

% ═══════════════════════════════════════════════════════════════
\section{Sujet 3 — La maison connectée}

\noindent\textbf{Examen :} BEPC \quad\textbf{Épreuve :} Physique-Chimie-Technologie
\quad\textbf{Durée :} 2 heures \quad\textbf{Coefficient :} 2 \quad\textbf{Notée sur :}
20 points.
\par\smallskip\noindent
L'épreuve comporte deux parties A et B, indépendantes et obligatoires.

\subsection*{Partie A — Évaluation des ressources (10 points)}

\paragraph{a) Évaluation des savoirs (5 points)}
\begin{enumerate}[label=\arabic*)]
\item Qu'est-ce que la valeur efficace d'une tension alternative ? \hfill (1pt)
\item Quel est le rôle d'un disjoncteur différentiel ? \hfill (1pt)
\item Citer deux fractions obtenues par distillation du pétrole brut. \hfill (1pt)
\item Qu'est-ce qu'un polymère ? \hfill (1pt)
\item Donner la relation entre la puissance électrique $P$, la tension $U$ et l'intensité
$I$. \hfill (1pt)
\end{enumerate}

\paragraph{b) Évaluation des savoir-faire (5 points)}
\begin{enumerate}[label=\arabic*)]
\item Un circuit domestique est protégé par un disjoncteur de calibre $I=20$ A sous une
tension $U=220$ V. Calculer la puissance maximale admissible $P_{\max}$ sur ce circuit.
\hfill (2pt)
\item Un appareil de puissance $P=1\,500$ W fonctionne pendant $t=3$ heures. Calculer
l'énergie consommée, en Wh puis en kWh, puis son coût au tarif de $79$ FCFA/kWh.
\hfill (2pt)
\item Sur un circuit protégé par un disjoncteur de calibre $16$ A sous $220$ V, on branche
trois appareils de $1\,200$ W, $800$ W et $900$ W. Le disjoncteur va-t-il se déclencher ?
Justifier. \hfill (1pt)
\end{enumerate}

\subsection*{Partie B — Évaluation des compétences (10 points)}

\noindent\textbf{Situation :} Une famille vient d'équiper sa maison de nouveaux appareils
électriques (four, téléviseur) et souhaite vérifier que son installation électrique peut
les supporter en toute sécurité. Elle s'interroge également sur le devenir des emballages
plastiques de ces nouveaux appareils.

\begin{enumerate}[label=\textbf{Tâche \arabic*.}, leftmargin=2.2cm]
\item La cuisine est protégée par un disjoncteur de calibre $I=16$ A sous une tension
$U=220$ V. Calculer la puissance maximale $P_{\max}$ admissible sur ce circuit.
\hfill (2pt)
\item La famille souhaite brancher simultanément, sur ce même circuit, un four de
$1\,800$ W et une bouilloire électrique de $1\,100$ W. Calculer la puissance totale
appelée et indiquer si le disjoncteur va se déclencher. \hfill (2pt)
\item Le nouveau téléviseur, de puissance $P=150$ W, fonctionne en moyenne $t=4$ heures
par jour. Calculer l'énergie qu'il consomme chaque jour, en Wh puis en kWh. \hfill (2pt)
\item Au tarif de $79$ FCFA/kWh, calculer le coût de fonctionnement du téléviseur sur un
mois de $30$ jours. \hfill (2pt)
\item Les emballages des nouveaux appareils portent un symbole triangulaire numéroté.
Expliquer à quoi sert ce symbole, puis proposer un geste permettant de réduire l'impact
environnemental de ces emballages plastiques. \hfill (2pt)
\end{enumerate}

% ═══════════════════════════════════════════════════════════════
\section{Sujet 3 — Corrigé}

\subsection*{Partie A}

\paragraph{a) Savoirs}
\begin{enumerate}[label=\arabic*)]
\item C'est la tension continue équivalente qui produirait le même effet (chauffage,
éclairage) que la tension alternative réelle.
\item Il protège les personnes : il détecte une différence entre l'intensité entrante et
sortante d'un circuit (fuite de courant) et coupe immédiatement l'alimentation.
\item Par exemple : l'essence et le gasoil (ou tout autre couple parmi gaz, kérosène,
fioul lourd, bitume).
\item Un polymère est une longue molécule formée par l'enchaînement répété de petites
molécules identiques appelées monomères.
\item $P=U\times I$.
\end{enumerate}

\paragraph{b) Savoir-faire}
\begin{enumerate}[label=\arabic*)]
\item $P_{\max}=U\times I=220\times20=4\,400$ W.
\item $E=P\times t=1\,500\times3=4\,500$ Wh $=4{,}5$ kWh. Coût $=4{,}5\times79=355{,}5$
FCFA.
\item $P_{\max}=220\times16=3\,520$ W. Puissance totale $=1\,200+800+900=2\,900$ W
$<3\,520$ W : le disjoncteur ne se déclenche pas.
\end{enumerate}

\subsection*{Partie B}

\textbf{Tâche 1.} $P_{\max}=U\times I=220\times16=3\,520$ W.

\textbf{Tâche 2.} Puissance totale $=1\,800+1\,100=2\,900$ W $<3\,520$ W : le disjoncteur
ne se déclenche pas, le branchement est possible.

\textbf{Tâche 3.} $E=P\times t=150\times4=600$ Wh $=0{,}6$ kWh par jour.

\textbf{Tâche 4.} Coût mensuel $=0{,}6\times79\times30=1\,422$ FCFA.

\textbf{Tâche 5.} Ce symbole indique le type de polymère utilisé et facilite le tri des
déchets en vue du recyclage. Pour réduire l'impact environnemental, on peut par exemple
trier systématiquement ces emballages afin qu'ils soient recyclés plutôt que jetés avec
les ordures ménagères.
